\section{Limitations, Failure Modes, and Non-Claims}
\label{sec:limitations}

\subsection{Evidence is first-party and retrospective}

The current live cases are maintained by the same ecosystem that developed
KFD. Public commits, immutable cuts, review records, and explicit falsifiers
make the claims inspectable but do not remove correlated incentives. The
Consequential Settlement mapping was reconstructed after action. Independent
adopters and prospective trials are required.

\subsection{Survivorship and narrative bias}

Once a candidate is named, many histories can be retold as if they led
inevitably to it. The mechanism counters this with frozen cuts, method
plurality, conservative baselines, false-candidate retention, and held-out
evaluation, but the current evidence does not yet demonstrate that those gates
are sufficient.

\subsection{Capture can distort action}

Rich Episode capture may change participant behavior, create surveillance
pressure, and privilege what is machine-recordable. A causal graph can still
exclude consequences that occur outside the capture boundary. Declared gaps do
not eliminate this harm. In sensitive settings, less retention or no capture
may be the correct choice.

\subsection{Ontology growth can export complexity}

A false Primitive replaces local reconstruction with global schema,
governance, migration, and training cost. Deletion and fuse tests are
necessary, not sufficient. Progressive disclosure and rival representations
must remain available. ``No new Primitive'' is a successful outcome.

\subsection{No automatic epistemic justice}

Perspective metadata does not guarantee that marginalized participants can
contribute, contest, or control promotion. The corpus owner may still dominate
access and interpretation. The governance requirements in
Section~\ref{sec:question-rights} are research and design obligations, not a
proof that a deployed institution is legitimate.

\subsection{Explicit non-claims}

This alpha paper does not claim:

\begin{itemize}[leftmargin=*]
  \item that KFD-6 is active or that a conforming autonomous loop exists;
  \item that every trace should become an Episode or every Episode be retained;
  \item that more data produces more Primitives;
  \item that perspective transformation dominates other genesis methods;
  \item that the six current candidate tracks are historically novel,
        universally minimal, or cross-domain requirements;
  \item that a schema validator proves causal truth, discovery, or promotion;
  \item that one physical object model is required; or
  \item that technical stewardship resolves legal, ethical, labor, cultural,
        or political authority.
\end{itemize}

These boundaries are part of the result. A theory of ontology revision that
cannot preserve the possibility that it is wrong would reproduce the failure
it claims to solve.
